\begin{abstract}
Optimizing an Apache Spark pipeline is not free: tuning execution parameters or compacting
small files itself consumes compute that must be recovered by future executions. We present
CAO-DF, a lightweight and interpretable cost-benefit framework that decides, for a given
workload and an expected number of future runs, whether to keep a static configuration, tune
the number of shuffle partitions, or compact the input files. The decision rests on an explicit
total-cost model and an optimization break-even point. We evaluate it in a controlled local
study of seeded synthetic workloads that vary data volume, file fragmentation and shuffle
volume, using a randomized-block protocol with 250 measured executions.
\tbd{A1 finding with numbers}. \tbd{A2 finding with numbers}. \tbd{Break-even finding}.
On \tbd{Z} unseen workloads the framework reduces estimated total cost by \tbd{Y\%} relative
to the static configuration while avoiding \tbd{W} unprofitable optimizations. Costs are
estimated from published cloud pricing, not billed, and all experiments run on a single machine.
\end{abstract}

\begin{IEEEkeywords}
Apache Spark, cost-aware optimization, break-even analysis, shuffle partitions,
file compaction, data pipelines, lakehouse, workload characterization
\end{IEEEkeywords}
